\begin{ouverture}[courbe]{3}{Multiples, diviseurs et nombres premiers}{Partie 1 · Nombres et calculs}
{Divisibilité · Décomposition · Fractions irréductibles · Périodicités}
\objectif{Décomposer un entier en facteurs premiers.}
\objectif{Simplifier une fraction et résoudre des partages.}
\objectif{Organiser une recherche de multiples communs.}
\prerequis{Division euclidienne et critères de divisibilité par 2, 3, 5 et 9.}
\end{ouverture}
\pagestyle{otMain}\essentiel\fontsize{11}{15}\selectfont
\essSection{1}{Divisibilité et nombres premiers}
\begin{definition}Un entier $a$ est un multiple d’un entier non nul $b$ lorsque $a=bk$ pour un entier $k$. Alors $b$ est un diviseur de $a$. Un nombre premier est un entier supérieur ou égal à 2 ayant exactement deux diviseurs positifs : 1 et lui-même.\end{definition}
\begin{exemple}2, 3, 5, 7, 11 et 13 sont premiers. Le nombre 1 ne l’est pas. Le nombre 91 ne l’est pas non plus : $91=7\times13$.\end{exemple}
\begin{propriete}Un entier est divisible par 2 si son dernier chiffre est pair, par 5 s’il se termine par 0 ou 5, par 3 ou par 9 si la somme de ses chiffres l’est. Un nombre peut vérifier plusieurs critères.\end{propriete}
\begin{savoirfaire}[Décomposer]Diviser successivement par des nombres premiers. Par exemple $60=2^2\times3\times5$ et $360=2^3\times3^2\times5$. Continuer jusqu’à ce que tous les facteurs soient premiers.\end{savoirfaire}
\begin{attention}Pour prouver qu’un nombre est composé, un seul diviseur autre que 1 et lui-même suffit. Pour prouver qu’il est premier, il faut exclure tous les diviseurs nécessaires, pas seulement 2 et 3.\end{attention}
\clearpage
\essSection{2}{Simplifier et organiser des groupes}
\begin{definition}Une fraction d’entiers est irréductible lorsque numérateur et dénominateur n’ont aucun diviseur commun positif autre que 1.\end{definition}
\begin{exemple}$84/150=(2^2\times3\times7)/(2\times3\times5^2)=14/25$. Les facteurs communs ont été simplifiés ; 14 et 25 n’en ont plus.\end{exemple}
\begin{savoirfaire}[Le plus grand nombre de groupes identiques]Le nombre de groupes doit diviser chacun des stocks. Chercher leur plus grand diviseur commun, à l’aide des décompositions ou de la liste des diviseurs.\end{savoirfaire}
\begin{savoirfaire}[La première rencontre de phénomènes périodiques]Le temps écoulé doit être un multiple de chaque période. Chercher le plus petit multiple commun strictement positif. Calendriers, engrenages et cycles d’éclosion d’espèces animales relèvent de ce raisonnement.\end{savoirfaire}
\begin{attention}Le mot « plus grand » ou « premier » demande une justification de l’optimalité. Donner un partage possible ou une rencontre possible ne suffit pas.\end{attention}
\otIntertitre{Avec un tableur}
Une colonne d’entiers et une colonne de restes permettent de tester les diviseurs. Contrôler la liste obtenue avec une décomposition en facteurs premiers.
\clearpage\rubrique{Découverte}
\begin{blocLivre}{ex-3-1}\begin{exercice}[1]{Tester des diviseurs}
Parmi 2, 3, 5, 9 et 10, lesquels divisent 2340 ?
\end{exercice}\end{blocLivre}
\begin{blocLivre}{ex-3-2}\begin{exercice}[1]{Premier ou composé ?}
Parmi 1, 2, 17, 21 et 49, relever les nombres premiers et justifier les exclusions.
\end{exercice}\end{blocLivre}
\begin{blocLivre}{ex-3-3}\begin{exercice}[1]{Décomposer un entier}
Décomposer 420 en facteurs premiers.
\end{exercice}\end{blocLivre}
\begin{blocLivre}{ex-3-4}\begin{exercice}[1]{Une fraction irréductible}
Rendre $198/252$ irréductible.
\end{exercice}\end{blocLivre}
\begin{blocLivre}{ex-3-5}\begin{exercice}[1]{Écrire des multiples}
Lister les multiples de 14 compris entre 70 et 130, bornes comprises.
\end{exercice}\end{blocLivre}

\clearpage\rubrique{Entraînement}
\begin{blocLivre}{ex-3-6}\begin{exercice}[2]{Tous les diviseurs}
Donner tous les diviseurs positifs de 72 et expliquer pourquoi la liste est complète.
\end{exercice}\end{blocLivre}
\begin{blocLivre}{ex-3-7}\begin{exercice}[2]{Un chiffre inconnu}
Déterminer les chiffres $a$ pour lesquels le nombre $45a$ est divisible par 3 et par 2.
\end{exercice}\end{blocLivre}
\begin{blocLivre}{ex-3-8}\begin{exercice}[2]{Une grande fraction}
Réduire $\dfrac{2^4\times3^2\times7}{2^2\times3^3\times5}$.
\end{exercice}\end{blocLivre}
\begin{blocLivre}{ex-3-9}\begin{exercice}[2]{Un carré parfait ?}
Le nombre $2^4\times3^3$ est-il un carré d’entier ? Par quel plus petit entier positif le multiplier pour obtenir un carré ?
\end{exercice}\end{blocLivre}
\begin{blocLivre}{ex-3-10}\begin{exercice}[2]{Un calcul par tableur}
Une colonne A contient les entiers de 1 à 20. La cellule B2 contient la formule \code{=MOD(84;A2)}. Que signifie un résultat nul ? Quels diviseurs de 84 inférieurs ou égaux à 20 obtient-on ?
\end{exercice}\end{blocLivre}

\clearpage\rubrique{Problèmes}
\begin{blocLivre}{ex-3-11}\begin{exercice}[3]{Préparer des lots}
On possède 156 cartes rouges et 204 cartes bleues. Former le plus grand nombre de lots identiques sans reste, avec les deux couleurs dans chaque lot. Donner le nombre de lots et leur contenu.
\end{exercice}\end{blocLivre}
\begin{blocLivre}{ex-3-12}\begin{exercice}[3]{Des cigales périodiques}
Dans un modèle simplifié, une espèce de cigales sort de terre tous les 13 ans et une autre tous les 17 ans. Les deux sont sorties la même année en 2024. En quelle année cette coïncidence se reproduira-t-elle pour la première fois ?
\end{exercice}\end{blocLivre}
\begin{blocLivre}{ex-3-13}\begin{exercice}[3]{Daller un rectangle}
Un rectangle de 180 cm sur 252 cm doit être recouvert de carrés identiques, sans découpe. Les côtés des carrés sont des nombres entiers de centimètres. Trouver le plus grand côté possible et le nombre de carrés.
\label{fig-ex-3-13}\figureGeo{dallage}
\end{exercice}\end{blocLivre}
\begin{blocLivre}{ex-3-14}\begin{exercice}[3]{Un stock à retrouver}
Un stock compte entre 500 et 600 objets. Il se range sans reste par paquets de 8, de 9 ou de 15. Déterminer toutes les possibilités.
\end{exercice}\end{blocLivre}
\begin{blocLivre}{ex-3-15}\begin{exercice}[3]{Des engrenages}
Deux roues de 24 et 40 dents engrènent. Deux dents marquées sont en contact au départ. Après combien de passages de dents retrouve-t-on ce contact ? Combien de tours fait chaque roue ?
\label{fig-ex-3-15}\figureGeo{engrenages}
\end{exercice}\end{blocLivre}

\clearpage
